\chapter{Pruebas}
\label{chap:pruebas}

\lettrine{L}{a} estrategia de pruebas de PadelCenter sigue la pirámide de
pruebas clásica: una base amplia de tests unitarios rápidos, una capa intermedia
de tests de integración que verifican el comportamiento con recursos reales, y
tests de arquitectura que garantizan el cumplimiento de las reglas de capas. Este
enfoque maximiza la confianza en el sistema minimizando el tiempo total de ejecución
del conjunto de pruebas. La escritura de tests unitarios antes o inmediatamente
después de cada caso de uso, descrita por Beck como práctica central del
desarrollo dirigido por pruebas~\cite{beck2002tdd}, ha guiado además el diseño
de las interfaces de los puertos del dominio: un puerto difícil de testear de
forma aislada es, en la mayoría de los casos, síntoma de un diseño con demasiadas
responsabilidades.

\section{Estrategia de pruebas}
\label{sec:estrategia_pruebas}

\begin{description}
  \item[Tests unitarios:] Verifican el comportamiento de un caso de uso o un
    componente de dominio de forma aislada, con las dependencias sustituidas por
    dobles de prueba (Mockito). Son muy rápidos y no requieren infraestructura
    externa.

  \item[Tests de integración:] Verifican el comportamiento del sistema de extremo
    a extremo a nivel HTTP, usando una base de datos PostgreSQL real gestionada
    por Testcontainers. Detectan errores que los tests unitarios no pueden encontrar:
    consultas JPA incorrectas, problemas de transaccionalidad, validaciones de
    esquema, etc.

  \item[Tests de arquitectura:] Verifican que la estructura del código cumple
    las reglas de la arquitectura hexagonal (ninguna clase del dominio importa
    clases de Spring, los casos de uso no dependen de JPA, etc.). Se ejecutan
    en cada build y fallan si alguien introduce una dependencia prohibida.
\end{description}

\section{Tests unitarios con Mockito}
\label{sec:tests_unitarios}

Los tests unitarios usan Mockito~\cite{mockito} y siguen el patrón \textit{given/when/then} y la convención
de nombres \texttt{should\{Behavior\}When\{Context\}()}. Por ejemplo:

\begin{verbatim}
@Test
void shouldThrowConflictExceptionWhenBookingOverlaps() {
    // given
    given(bookingRepository.existsConflict(FIELD_ID, START, END))
        .willReturn(true);

    // when / then
    assertThatThrownBy(() -> useCase.execute(command))
        .isInstanceOf(BookingConflictException.class);
}
\end{verbatim}

Cada clase de caso de uso tiene su correspondiente clase de test. Los dobles
de prueba se crean con \texttt{@Mock} y la clase bajo test con \texttt{@InjectMocks}
en combinación con \texttt{MockitoExtension}.

Los objetos de dominio que contienen lógica de negocio (por ejemplo, la
validación de solapamiento de horarios en \texttt{Booking}) también tienen tests
unitarios propios, sin mocks, que verifican el comportamiento con entradas y
salidas concretas.

\section{Tests de integración con Testcontainers}
\label{sec:tests_integracion}

Los tests de integración usan \texttt{@SpringBootTest(webEnvironment = RANDOM\_PORT)}
y un contenedor PostgreSQL gestionado por Testcontainers~\cite{testcontainers}:

\begin{verbatim}
@Container
static PostgreSQLContainer<?> postgres =
    new PostgreSQLContainer<>("postgres:16")
        .withDatabaseName("padelcenter_test");
\end{verbatim}

El contenedor se levanta una sola vez por clase de test
(\texttt{@TestcontainersExtension}) para minimizar el overhead. Flyway aplica
las migraciones automáticamente al arrancar el contexto Spring, garantizando que
el esquema de test siempre coincide con el de producción.

Las peticiones HTTP se realizan con \texttt{TestRestTemplate} o con el cliente
Orval generado en modo test, cubriendo el ciclo completo:

\begin{enumerate}
  \item Autenticación previa para obtener un JWT válido.
  \item Llamada al endpoint bajo prueba con el token en el header.
  \item Verificación del código de estado HTTP y del cuerpo de la respuesta.
  \item Verificación directa en base de datos para casos críticos de escritura.
\end{enumerate}

\section{Tests de arquitectura con ArchUnit}
\label{sec:tests_arquitectura}

ArchUnit~\cite{archunit} permite expresar las reglas de arquitectura como código Java y
verificarlas en tiempo de compilación. PadelCenter define las siguientes reglas:

\begin{itemize}
  \item Las clases del paquete \texttt{domain} no deben depender de ninguna
    clase de los paquetes \texttt{org.springframework} ni
    \texttt{jakarta.persistence}.
  \item Las clases del paquete \texttt{application} solo pueden depender de
    \texttt{domain}, no de \texttt{infrastructure}.
  \item Todas las clases de casos de uso deben tener exactamente un método
    público llamado \texttt{execute}.
  \item Los nombres de las clases de repositorio JPA deben terminar en
    \texttt{JpaRepository}.
\end{itemize}

Si alguna de estas reglas se viola, el test falla con un mensaje descriptivo que
indica exactamente qué clase incumple la regla y por qué.

\section{Resumen de cobertura}
\label{sec:resumen_cobertura}

La tabla~\ref{tab:resumen_tests} muestra el número de tests por tipo y la
cobertura global del backend, medida con JaCoCo~\cite{jacoco} a partir de la
fusión de los datos de ejecución de Surefire (unitarios) y Failsafe
(integración) — \texttt{mvn clean verify}. El código generado a partir del
contrato OpenAPI (DTOs e interfaces de controlador) se excluye del cálculo,
ya que no es código escrito a mano.

\begin{table}[h]
\centering
\resizebox{\textwidth}{!}{%
\begin{tabular}{lccc}
\hline
\textbf{Tipo} & \textbf{Tests} & \textbf{Capas cubiertas} & \textbf{Cobertura (líneas)} \\
\hline
Unitarios (Mockito)          & 193 & \texttt{domain}, \texttt{application} & 85,2\% \\
Integración (Testcontainers) & 35  & \texttt{infrastructure} + end-to-end   & 49,4\% \\
Arquitectura (ArchUnit)      & 5   & Reglas de capas                        & N/A    \\
\hline
\textbf{Total}               & \textbf{233} & \textbf{---} & \textbf{64,3\%} \\
\hline
\end{tabular}}
\caption{Resumen de tests y cobertura del backend (medida con JaCoCo, perfil \texttt{mvn clean verify}).}
\label{tab:resumen_tests}
\end{table}

La cobertura del 85,2\% en \texttt{domain} y \texttt{application} refleja que
las capas más críticas del sistema están ampliamente verificadas mediante
tests unitarios rápidos. La capa de infraestructura queda más rezagada
(49,4\%): además de los adaptadores de difícil cobertura unitaria
(configuración de Spring Security, mapeo JPA), faltan clases
\texttt{*ControllerIT} dedicadas para los recursos de torneos, partidos y
autenticación — actualmente solo existen para usuarios, centros, pistas y
reservas (\texttt{UserControllerIT}, \texttt{CenterControllerIT},
\texttt{FieldControllerIT}, \texttt{BookingControllerIT}) y para la
configuración de seguridad (\texttt{SecurityConfigPublicEndpointsIT}). Esta
brecha queda identificada como trabajo futuro en la
sección~\ref{chap:conclusions}.
